% 05/11/2002 at 10:00  
% Din�mica do Ponto Material: A Lei de Hooke
% Written by TORDAR 

\documentclass[12pt]{article}  

\pagestyle{empty}

\usepackage{graphicx} 
\usepackage{pst-all} 

%Let's go: 

\begin{document} 
\null\vskip4cm
\begin{center} 
\begin{pspicture}(-0.1,-0.1)(1.1,1.1)
\psset{unit=8cm}
\psline[linewidth=2pt,linecolor=black,linestyle=dashed,dash=5mm 1mm](0.25,0)(0.75,0)
\psline[linewidth=2pt,linecolor=black,linestyle=dashed,dash=5mm 1mm](0.4,0.14)(0.75,0.14)
\psframe[linewidth=1pt,
         linecolor=lightgray,
	 dimen=inner,
	 fillstyle=hlines*,
	 fillcolor=lightgray,
	 hatchangle=45](0.1,0.7)(0.4,0.8)
\psframe[linewidth=1pt,
         linecolor=lightgray,
	 dimen=inner,
	 fillstyle=hlines*,
	 fillcolor=lightgray,
	 hatchangle=45](0.6,0.7)(0.9,0.8)
\pscoil[linewidth=2pt,linecolor=blue,coilarm=0.5cm,coilwidth=0.5cm]{|-|}(0.25,0.7)(0.25,0)
\scalebox{1 0.8}{
\pscoil[linewidth=2pt,linecolor=blue,coilarm=0.5cm,coilwidth=0.5cm]{|-|}(0.75,0.875)(0.75,0.175)
                 }
\pcline[linewidth=2pt,linecolor=black,offset=15pt]{<->}(0.1,0)(0.1,0.7)
\lput*{:U}{\scalebox{1.5}{$l$}}
\pcline[linewidth=5pt,linecolor=black,offset=15pt]{->}(0.73,0)(0.73,0.14)
\pcline[linewidth=5pt,linecolor=blue]{<-}(0.73,0)(0.73,0.14)
\rput{0}(0.57,0.07){\scalebox{1.5}{$\Delta\mathbf{r}$}}
\rput{0}(0.83,0.07){\scalebox{1.5}{$\mathbf{F}_m$}}
\end{pspicture}
\end{center}
\end{document} 
